\chapter{Structure from Constraint}
\label{ch:structure-from-constraint}

Before any angle is measured or any triangle solved, it is worth asking a more
primitive question: what does it mean for a mathematical object to have
structure at all? This short opening chapter proposes an answer using discrete
state spaces, in which the objects of study are finite collections of symbols,
codes, or configurations rather than continuous geometric figures. The purpose
is not to anticipate the content of later chapters but to introduce, in the
simplest possible setting, the two ideas that organize the entire book:
constraint and invariance.

\section{Discrete State Spaces}

Consider a finite set of symbols, called an \emph{alphabet}, and the collection
of all finite sequences that can be formed from it. If the alphabet has $k$
symbols, then there are $k^n$ sequences of length $n$, and this count grows
without any internal organization worth remarking on: every sequence is as
good as any other. Structure enters only once some sequences are declared
admissible and others are not. It is the constraint, not the alphabet, that
produces something worth studying.

\begin{definition}
A \emph{discrete state space} is a pair $(\Sigma, n)$ consisting of a finite
alphabet $\Sigma$ and a length $n$, together with the set $\Sigma^n$ of all
length-$n$ sequences over $\Sigma$. An \emph{admissibility constraint} on
$(\Sigma,n)$ is a rule that partitions $\Sigma^n$ into two classes: the
\emph{admissible} sequences, which satisfy the rule, and the
\emph{inadmissible} sequences, which do not. The set of admissible sequences
is written $\mathcal{A} \subseteq \Sigma^n$.
\end{definition}

A constraint of this kind rarely acts on a sequence all at once. Far more
commonly it acts \emph{locally}: whether a given symbol may follow the
previous one depends only on that previous symbol, not on the entire history
of the sequence. Such constraints are conveniently recorded in a
\emph{transition graph}: a directed graph whose vertices are the symbols of
$\Sigma$, with an edge from symbol $a$ to symbol $b$ exactly when $b$ is
permitted to follow $a$. This single graph then governs every sequence in
$\mathcal{A}$, regardless of length.

\section{Constraints Generate Structure}

\begin{theorem}
Let $(\Sigma, n)$ be a discrete state space whose admissibility constraint is
local, recorded by a transition graph $G$ on vertex set $\Sigma$. Then the
admissible sequences of length $n$ are in bijection with the directed walks of
length $n-1$ in $G$. Consequently, if $M$ is the adjacency matrix of $G$, the
number of admissible sequences of length $n$ beginning at symbol $a$ and
ending at symbol $b$ is the $(a,b)$ entry of $M^{n-1}$.
\end{theorem}

\begin{proof}
A sequence $s_1 s_2 \cdots s_n \in \Sigma^n$ is admissible exactly when each
consecutive pair $(s_i, s_{i+1})$ corresponds to an edge of $G$, for
$i = 1, \dots, n-1$. This is precisely the condition for $s_1, s_2, \dots,
s_n$ to trace a directed walk of length $n-1$ through $G$, visiting vertices
$s_1$ through $s_n$ in order. The correspondence between admissible sequences
and such walks is one-to-one by construction. The entry-counting claim then
follows from the standard fact that the $(a,b)$ entry of $M^{n-1}$ counts
directed walks of length $n-1$ from $a$ to $b$, which is proved by induction
on $n$ using the definition of matrix multiplication.
\end{proof}

What the theorem makes precise is something one can already sense
informally: once a constraint is imposed, the admissible set stops being an
undifferentiated pile of sequences and becomes something with countable,
predictable structure, generated by iterating a single local rule. This is
the first appearance, in the simplest setting available, of a pattern that
will recur throughout the book. A constraint restricts a large space of
possibilities to a smaller admissible subset, and that subset, far from being
arbitrary, inherits a structure directly traceable to the constraint that
produced it.

\begin{example}
Let $\Sigma = \{0, 1, \dots, 9\}$ and consider the discrete state space of
length-$n$ digit strings representing $n$-digit whole numbers, with the
constraint that the first symbol $s_1$ may not be $0$. This is a local
constraint concentrated entirely on the first position: every symbol may
follow every other symbol in positions $2$ through $n$, but only the symbols
$1$ through $9$ are admissible in position $1$. The number of admissible
sequences is therefore $9 \cdot 10^{n-1}$, obtained by choosing among nine
possibilities for the first digit and ten possibilities for each remaining
digit. This count is the familiar number of $n$-digit whole numbers, and it
follows immediately from treating positional notation as a constrained
discrete state space rather than as a definition to be taken for granted.
\end{example}

\begin{practicenow}
Let $\Sigma = \{0,1\}$ and consider binary strings of length $n$ with the
constraint that no two consecutive symbols may both equal $1$. Draw the
transition graph on $\{0,1\}$ implied by this constraint, and use it to count
the admissible strings of length $4$ by hand. Compare the result to the
count predicted by the adjacency-matrix method of the theorem.
\end{practicenow}

\begin{sidebar}{Cistercian Numerals and Independent Composition}
A striking illustration of composition from constrained parts comes from
the Cistercian numeral system, developed by medieval monks to write any
whole number from one to nine thousand nine hundred ninety-nine as a
single glyph. A vertical stem carries up to four attached strokes, one in
each of four fixed positions corresponding to units, tens, hundreds, and
thousands, and each position independently carries one of ten possible
strokes, drawn from a shared base pattern rotated and reflected into the
position it occupies. The system is therefore generated by choosing one
symbol from a set of ten for each of four independent positions, giving
$10^4 = 10{,}000$ representable values once the empty choice at every
position, representing zero, is included. Nothing about the choice made in
the units position constrains the choice available in the tens position;
the four positions act as independent coordinates of a single glyph, in
exactly the sense that a vector in this book is built from independent
coordinates along different axes. What makes such a system powerful is
not the alphabet of ten strokes by itself, but the fact that composing
strokes across four independent positions multiplies the number of
representable values far beyond what any single position could express
alone, a pattern that will reappear once vectors and matrices are built
from independently varying components of their own.
\end{sidebar}

\section{From Discrete Constraint to Continuous Invariance}

The state spaces considered so far have been finite, and their constraints
have been combinatorial rules about which symbols may follow which. Nothing
in the underlying logic, however, depends on finiteness. A continuous
analogue appears as soon as one asks which points $(x,y)$ in the plane
satisfy $x^2 + y^2 = 1$. Here the ambient space $\mathbb{R}^2$ plays the role
of $\Sigma^n$, the equation plays the role of the admissibility constraint,
and the resulting unit circle plays the role of $\mathcal{A}$: a
lower-dimensional admissible set carved out of a larger, unconstrained space
by a single governing condition. Chapter~4 will take up this continuous case
in detail, but it is the same idea encountered here in miniature. A
constraint isolates an admissible substructure, and that substructure,
rather than the ambient space that contained it, becomes the proper object of
study. Much of what follows in this book can be read as a sustained
exploration of what happens when that single observation is taken seriously.

\section{Exercises}
\begin{exercise}
Let $\Sigma = \{a, b, c\}$ with transition graph given by the edges
$a \to b$, $b \to c$, $c \to a$, and $c \to b$. List all admissible sequences
of length $3$ beginning at $a$, and verify your count against the
adjacency-matrix method of the theorem.
\end{exercise}

\begin{exercise}
Explain, in a short paragraph, why the count $9 \cdot 10^{n-1}$ from the
worked example would change if the constraint instead forbade the digit $0$
from appearing in \emph{any} position rather than only the first. Compute the
new count and identify the transition graph that generates it.
\end{exercise}
